\documentclass[11pt]{article}
\usepackage[margin=1in]{geometry}
\usepackage{amsmath,amssymb,amsthm}
\usepackage{hyperref}
\usepackage{xurl}

\newtheorem{theorem}{Theorem}
\newtheorem{lemma}[theorem]{Lemma}
\newtheorem{corollary}[theorem]{Corollary}
\theoremstyle{definition}
\newtheorem{definition}[theorem]{Definition}
\theoremstyle{remark}
\newtheorem{remark}[theorem]{Remark}

\let\oldus\_
\renewcommand{\_}{\oldus\allowbreak}

\newcommand{\tr}{\operatorname{tr}}
\newcommand{\Tr}{\operatorname{Tr}}
\newcommand{\C}{\mathbb{C}}

\title{A proof of Conjecture 00000004780\\[2pt]
\large Swapping the diagonal blocks keeps every joint moment but changes the
$M_2(\C)$-valued distribution}
\author{Jackmeson1}
\date{}

\begin{document}
\sloppy
\maketitle

\section{The conjecture}

Conjecture 00000004780 reads (English version; the Chinese version says the same, with
``all joint moments asymptotically identical''):

\begin{quote}
\emph{Definition:} Joint moments and operator-valued distributions are two layers of
noncommutative invariants. \emph{Conjecture:} There exist two matrix models whose joint moments
have identical asymptotics while their operator-valued distributions differ, and the separation is
realized by an explicit rearrangement of the block structure.
\end{quote}

It is an existence statement. We give two explicit matrix models, the second obtained from the
first by swapping its two diagonal blocks, whose joint moments coincide for every matrix size and
every noncommutative polynomial, while their $M_2(\C)$-valued distributions differ for every matrix
size and in the limit. Everything below is proved in Lean 4 with Mathlib.

\section{Definitions}

Throughout, $N$ is the size parameter and matrices are written in block form
\[ M_2(M_N(\C)) = M_{2N}(\C), \qquad \text{indexed by } \{0,1\}\times\{0,\dots,N-1\}. \]
In Lean this is \texttt{BMat N := Matrix (Fin 2 $\times$ Fin N) (Fin 2 $\times$ Fin N) $\mathbb C$}.
The definitions below are stated for $N\ge1$. Lean also allows $N=0$ as a formal extension (the
matrices are empty, $b\mapsto b\otimes 1_0$ is not injective and $E_0=0$); this level plays no
role: every separating statement assumes $N\ge1$, and it does not affect any limit.

\begin{definition}[Matrix model, joint moments]
A matrix model of $r$ matrices is a sequence $X=(X_N)_{N}$ of $r$-tuples
$X_N=(X_{N,1},\dots,X_{N,r})$ in $M_2(M_N(\C))$ (Lean: \texttt{Model r}). Its joint moments are
\[ \varphi_N(p) = \tr_{2N}\bigl(p(X_{N,1},\dots,X_{N,r})\bigr),\qquad
\tr_{2N}=\tfrac{1}{2N}\Tr, \]
for every noncommutative polynomial $p\in\C\langle x_1,\dots,x_r\rangle$ (Lean:
\texttt{jointMoment X N p}, with $p$ in \texttt{FreeAlgebra $\mathbb C$ (Fin r)} and evaluation by
\texttt{FreeAlgebra.lift}). Taking $p$ a monomial gives the moments
$\tr(X_{i_1}\cdots X_{i_m})$; all matrices used below are self-adjoint, so $*$-moments are moments.
\end{definition}

\begin{definition}[Operator-valued distribution over $B=M_2(\C)$]
Let $B=M_2(\C)$ act on $M_2(M_N(\C))$ by $b\mapsto b\otimes 1_N$ (Lean: \texttt{emb}), and let
\[ E_N=\mathrm{id}_2\otimes\tr_N\colon M_2(M_N(\C))\to M_2(\C),\qquad
E_N(M)_{ij}=\frac1N\sum_{a} M_{(i,a),(j,a)} \]
be the block-wise normalized partial trace (Lean: \texttt{condExp}). For $N\ge1$ it is a unital,
$B$-bimodular projection onto $B\otimes 1_N$, i.e.\ the usual $B$-valued conditional expectation
for $2\times2$ block matrices.
The $B$-valued distribution of $X_N$ is the family of $B$-valued moments
\[ E_N\bigl(b_0\,X_{N,i_1}\,b_1\,X_{N,i_2}\cdots X_{N,i_m}\,b_m\bigr),\qquad b_k\in B \]
(Lean: \texttt{bMoment X N b0 l} with $l=[(i_1,b_1),\dots,(i_m,b_m)]$).
\end{definition}

\begin{definition}[Rearrangement of the block structure]
For a permutation $\sigma$ of the two block indices, let $P_\sigma=\sigma\otimes 1_N$ be the
permutation of $\{0,1\}\times\{0,\dots,N-1\}$ (Lean: \texttt{blockPerm}). The rearranged model is
$(\mathrm{rearrange}\ \sigma\ X)_{N,i}=P_\sigma X_{N,i}P_\sigma^{-1}$, which moves block $(k,l)$ of
$X_{N,i}$ to position $(\sigma k,\sigma l)$ (Lean: \texttt{rearrange}, via \texttt{Matrix.reindex}).
\end{definition}

\section{General facts about rearrangements}

\begin{lemma}[\texttt{rearrange\_jointMoment}]
For every model $X$, every $\sigma$, every $N$ and every $p$:
$\varphi_N^{\mathrm{rearrange}\,\sigma X}(p)=\varphi_N^{X}(p)$.
\end{lemma}

\begin{proof}
$M\mapsto P_\sigma M P_\sigma^{-1}$ is an algebra automorphism of $M_{2N}(\C)$
(\texttt{Matrix.reindexAlgEquiv}), so by the universal property of the free algebra
$p(P_\sigma X P_\sigma^{-1})=P_\sigma\,p(X)\,P_\sigma^{-1}$ (\texttt{eval\_reindex}), and the trace is
invariant under reindexing (\texttt{trace\_reindex\_self}).
\end{proof}

\begin{lemma}[\texttt{rearrange\_condExp}, \texttt{rearrange\_bMoment\_first}]
$E_N(P_\sigma M P_\sigma^{-1})=\sigma\,E_N(M)\,\sigma^{-1}$ for every $M$. Consequently, for every
model, $E_N\bigl((\mathrm{rearrange}\ \sigma X)_{N,i}\bigr)=\sigma\,E_N(X_{N,i})\,\sigma^{-1}$.
\end{lemma}

So a rearrangement never changes joint moments. On the full $B$-valued distribution it acts by
transport of structure: with $\alpha(b)=\sigma b\sigma^{-1}$, the moment of the rearranged model
with coefficients $b_0,\dots,b_m$ equals $\alpha$ applied to the moment of $X$ with coefficients
$\alpha^{-1}(b_0),\dots,\alpha^{-1}(b_m)$ (conjugating the output alone, with the coefficients
fixed, would be wrong in general). Lean proves only the special case with identity coefficients
(Lemma~2), which is all the example below needs. The explicit pair shows that the action can be
nontrivial.

\section{The explicit pair}

Let
\[ A=\begin{pmatrix}1&0\\0&0\end{pmatrix},\qquad C=\begin{pmatrix}0&1\\1&0\end{pmatrix},\qquad
X_N=(A\otimes 1_N,\ C\otimes 1_N),\qquad Y=\mathrm{rearrange}\,(0\ 1)\,X . \]
Thus $X_{N,1}=\begin{pmatrix}1_N&0\\0&0\end{pmatrix}$, $X_{N,2}=\begin{pmatrix}0&1_N\\1_N&0\end{pmatrix}$,
and $Y$ swaps the two diagonal blocks:
$Y_{N,1}=\begin{pmatrix}0&0\\0&1_N\end{pmatrix}$, $Y_{N,2}=X_{N,2}$.
All four matrices are self-adjoint (\texttt{X\_isHermitian}, \texttt{Y\_isHermitian}).
Lean indexes the two matrices by \texttt{Fin 2}, so $X_{N,1}$ is \texttt{X N 0} and the
$B$-valued moment $E_N(X_{N,1})$ is \texttt{bMoment X N 1 [(0, 1)]}.

\begin{theorem}[\texttt{conjecture4780}]
\begin{enumerate}
\item $Y=\mathrm{rearrange}\,(0\ 1)\,X$, and all matrices of both models are self-adjoint.
\item For every $N$ and every $p\in\C\langle x_1,x_2\rangle$: $\varphi_N^X(p)=\varphi_N^Y(p)$.
Both sequences converge, to the same limit $\tr_2(p(A,C))$.
\item For every $N\ge1$: $E_N(X_{N,1})=\mathrm{diag}(1,0)$ and $E_N(Y_{N,1})=\mathrm{diag}(0,1)$.
Hence the $B$-valued distributions of $X_N$ and $Y_N$ differ for every $N\ge1$.
\item The limits of these first $B$-valued moments are $\mathrm{diag}(1,0)\ne\mathrm{diag}(0,1)$,
so the limiting $B$-valued distributions differ as well. (Lean proves the convergence of these
first $B$-valued moments; see the remark on full limiting distributions below.)
\end{enumerate}
\end{theorem}

\begin{proof}
(1) holds by definition and because $A,C$ and $1_N$ are self-adjoint. (2) The equality is
Lemma 1. For the limit: $b\mapsto b\otimes 1_N$ is a unital algebra homomorphism
(\texttt{embHom}), so $p(X_N)=p(A,C)\otimes 1_N$, and
$\tr_{2N}(M\otimes 1_N)=\frac{\Tr M\cdot N}{2N}=\tr_2(M)$ for $N\ge1$ (\texttt{ntr\_emb},
\texttt{jointMoment\_X}); so the sequence is eventually constant. (3) The $B$-valued moment with
$b_0=b_1=1$ and the single letter $x_1$ is $E_N(X_{N,1})=E_N(A\otimes 1_N)=A$
(\texttt{condExp\_emb}); by Lemma 2 the same moment of $Y$ is $\sigma A\sigma^{-1}=\mathrm{diag}(0,1)$.
(4) Both sequences are eventually constant.
\end{proof}

\section{Remarks on the reading}

\begin{itemize}
\item \emph{Identical asymptotics.} We prove equality of the joint moments for every $N$, which
gives identical asymptotics under any reading (same limits, same rates, same expansions), and we
also prove that the limits exist.
\item \emph{Operator-valued distribution.} We use the standard $M_2(\C)$-valued distribution of a
$2\times2$ block matrix model with respect to $E_N=\mathrm{id}_2\otimes\tr_N$. The two models
already differ at the first $B$-valued moment, which belongs to every version of the $B$-valued
distribution (with or without $B$-coefficients, finite $N$ or limit).
\item \emph{Full limiting distributions.} The Lean theorem states convergence only for the first
$B$-valued moments. The full limiting distributions do exist: for $N\ge1$, multiplicativity of
$b\mapsto b\otimes1_N$ and $E_N(b\otimes1_N)=b$ give
$E_N(b_0X_{N,i_1}b_1\cdots X_{N,i_m}b_m)=b_0g_{i_1}b_1\cdots g_{i_m}b_m$ independently of $N$, with
$(g_1,g_2)=(A,C)$ for $X$ and $(\sigma A\sigma^{-1},C)$ for $Y$. This observation is not formalized;
the separation itself only needs the first moment.
\item \emph{Rearrangement.} The second model is literally the first with its block indices
permuted. Lemma 1 holds for every model and every $\sigma$, in particular for every sample of a
random matrix model.
\item \emph{Degenerate models.} The explicit models are deterministic (constant random matrices)
and have the form $a\otimes 1_N$. The conjecture does not require randomness or a particular
ensemble.
\end{itemize}

\section{Lean formalization and axiom audit}

The file \texttt{lean/Conjecture4780/Basic.lean} (248 lines, Lean 4.33.1, Mathlib v4.33.1)
contains \texttt{ntr}, \texttt{Model}, \texttt{eval}, \texttt{jointMoment}, \texttt{condExp},
\texttt{emb}, \texttt{bMoment}, \texttt{blockPerm}, \texttt{rearrange}, the lemmas
\texttt{rearrange\_jointMoment}, \texttt{rearrange\_condExp}, \texttt{rearrange\_bMoment\_first},
\texttt{condExp\_emb}, \texttt{ntr\_emb}, \texttt{jointMoment\_X}, and the main theorem
\texttt{Conjecture4780.conjecture4780}. It builds with \texttt{lake build}. For each of
\texttt{conjecture4780}, \texttt{rearrange\_jointMoment}, \texttt{rearrange\_condExp} and
\texttt{rearrange\_bMoment\_first}, \texttt{\#print axioms} reports exactly
\texttt{[propext, Classical.choice, Quot.sound]}. There is no \texttt{sorry}, no
\texttt{native\_decide} and no added axiom.

\end{document}
